% !PRINTED 475
\weberpage{475}
%<<CONT
nary linear substitutions for the modulus 2, which we investigated in \S.\,79\footnote{According to an oral communication, this is the way in which \gesperrt{Kronecker} formed them.}.

For this purpose we have to denote the seven quantities by three indices $(x_1, x_2, x_3)$, which are taken with respect to the modulus 2 and where the combination $(0, 0, 0)$ is excluded. One thus obtains the seven quantities

\weq{1}{(1,0,0),\ (0,1,0),\ (0,0,1);\ (0,1,1),\ (1,0,1),\ (1,1,0),\ (1,1,1),}

and if one replaces the $x_1, x_2, x_3$ by the linear combinations

\weq{2}{a x_1 + b x_2 + c x_3,\ a_1 x_1 + b_1 x_2 + c_1 x_3,\ a_2 x_1 + b_2 x_2 + c_2 x_3}

in which the substitution

\weq{3}{\begin{pmatrix} a, & b, & c \\ a_1, & b_1, & c_1 \\ a_2, & b_2, & c_2 \end{pmatrix}}

is to be taken with respect to the modulus 2, one obtains the permutation group $P_{168}$.

If, as in \S.\,79, we choose for $\tau$ the substitution of the congruence group

\weq{4}{\tau = \begin{pmatrix} 1, & 0, & 1 \\ 1, & 0, & 0 \\ 0, & 1, & 0 \end{pmatrix}}

then we can denote the quantities $z$ in such a way that they pass cyclically into one another by the application of $\tau$ and its repetitions, where we may choose an arbitrary one of the quantities (1) for $z_0$, say thus:

\weq{5}{z_0 = (1,0,0),\ z_1 = (1,1,0),\ z_2 = (1,1,1),\ z_3 = (0,1,1),\ z_4 = (1,0,1),\ z_5 = (0,1,0),\ z_6 = (0,0,1).}

Then $\omega$ is to be chosen so that $z_0, z_4, z_6$ remain unchanged by it and $z_1$ is interchanged with $z_2$, and $z_3$ with $z_5$. This gives

\weq{6}{\omega = \begin{pmatrix} 1, & 0, & 0 \\ 0, & 1, & 0 \\ 0, & 1, & 1 \end{pmatrix},}

and from this one obtains, by \S.\,72, (17):

\weq{7}{\chi = \begin{pmatrix} 1, & 0, & 0 \\ 0, & 0, & 1 \\ 0, & 1, & 1 \end{pmatrix}, \quad \Theta = \begin{pmatrix} 1, & 1, & 1 \\ 0, & 1, & 0 \\ 0, & 1, & 1 \end{pmatrix}.}
